%%%%%%%%%%%%%%%%%%%%%%%%%%%%%%%%%%%%%%%%%%%%%%%%%%%%%%%%%%%%%%%%%%%%%%%%%%%%%%%%%%%%%%%%%%%%%%%%%%%%%%%%%

%%Modelo desenvolvido por Alunos do IFSP - PEP.

%Autores: Christopher Alves de Oliveira, Inácio Ribeiro, Paulo Ricardo Servilha Sparapan.

%Professores: Alexandre Ataide Carniato e Diego Nunes da Silva

%Data: 18/08/2020.

%%%%%%%%%%%%%%%%%%%%%%%%%%%%%%%%%%%%%%%%%%%%%%%%%%%%%%%%%%%%%%%%%%%%%%%%%%%%%%%%%%%%%%%%%%%%%%%%%%%%%%%%%

    
\documentclass[a4paper,10pt]{article} %Selecionando a classe que gera artigos.
    
\input{Arquivos_Config/preambulo} %Pacotes principais


%%%%%%%%%%%%%%%%%%%%%%%%%%%%%%%%%%%%%%%%%%%%%%%%%%%%%%%%%%%%%%%%%%%%%%%%%%%%%%%%%%%%%%%%%%%%%%%%

\pagestyle{empty}


\begin{document} %Begin Inicia o Documento
 
\input{Arquivos_Config/cabecalho} %Alterar Título, autores, e apagar tabela !

\begin{center}
\textbf{TÍTULO(Helvetica, 9, Negrito, Centralizado)}\vspace{0.5cm}
    
FULANO C. SILVA (IFPB, Campus Sousa), AUTOR 2 (IFPB, Campus Cabedelo), AUTOR 3 (Centro de Educação, UFPB)
\\\color{red}OS NOMES DEVEM SER OMITIDOS PARA A SUBMISSÃO E ENTRAM APENAS NA VERSÃO FINAL\color{black}
    
    (Helvetica, 9, Centralizado, Máximo seis autores)\vspace{0.1cm}
\end{center}
\begingroup
 %   \fontsize{9pt}{11pt}\selectfont
  
 % $^1$Graduando em Tecnologia de Análise e Desenvolvimento de Sistemas, Bolsista PIBIFSP, IFSP, Câmpus Presidente Epitácio, fulanocsilva@ifsp.edu.br. (Helvetica, 9, Justificado)
  
 % $^2$
  
%  $^3$
  
%  $^4$
\small 
\noindent \textbf{E-mails:} fulano@academico.ifpb.edu.br, autor2@ifpb.edu.br, autor3@ufpb.br.
\\
\noindent\textbf{Área de conhecimento:(Tabela CNPq)}: 1.03.03.04-9 Sistemas de Informação. 
\\
\noindent\textbf{Palavras-Chave}: máximo de seis, separadas por ponto e vírgula (;), procurando não repetir palavras do título, escritas em letras minúsculas. (Helvetica, 8, Justificado).
\endgroup

%----------------------------------------TEXTOS-------------------------------------------------%

\section{Introdução}

O propósito destas instruções é orientar aos autor(es) quanto à formatação dos resumos expandidos a serem submetidos ao SIMPIF 2023 do Instituto Federal de Educação, Ciência e Tecnologia da Paraíba. O texto com as instruções e em parênteses devem ser removidos do documento final. Além disso, o documento final deve ter até 03 páginas, desconsiderando o espaço reservado aos nomes dos autores.

As referências devem ser grafadas no final do resumo, em ordem alfabética e cronológica. Só devem compor as Referências as fontes que tenham sido efetivamente citadas ao longo do texto. Exemplos de referência no texto: \cite{galvani}, \cite{national_instruments}.

A introdução deve descrever o estudo de forma clara e objetiva, destacando a relevância do problema investigado, com base na literatura, e os principais objetivos do trabalho. (Helvetica, 9, Justificado, Máximo 20 linhas).

\section{Materiais e Métodos}

Descrever, de forma objetiva, sobre como o trabalho foi realizado. (Helvetica, 9, Justificado).

\subsection{Modelo de Equação}

A Equação~\ref{eq:nome} calcula o IC.

\begin{equation} \label{eq:nome}
IC=\frac{F*9,81}{A}*10^{-6},
\end{equation}
em que, $IC$ é o índice de cone em $MPa$, $F$ é a força em $kgf$ e $A$ é a área do cone em $m^{2}$.

\section{Resultados e Discussão}

Informar e/ou comentar a respeito dos resultados obtidos com a pesquisa.

\subsection{Gráficos}

Devem apresentar-se sem bordas, descritos com o mesmo tipo e tamanho de letras contidas no texto e a legenda na posição inferior do mesmo. A numeração deve ser sucessiva em algarismos arábicos.

\subsection{Tabelas} 

Evitar tabelas extensas e dados supérfluos; adequar seus tamanhos ao espaço útil do papel e colocar, na medida do possível, apenas linhas contínuas horizontais; suas legendas devem ser concisas e autoexplicativas. Na discussão, confrontar os dados obtidos com a literatura. (Helvetica, 9, Justificado).

\subsection{Modelo de Figuras}

A Figura~\ref{fig:nome_referencia_figura1} mostra a logo do IFPB.

\begin{figure}[h!]
\centering
\includegraphics[width=3cm,angle=0]{Figuras/logo_ifpb_vert.png}
\caption{Logo do IFPB.}
\label{fig:nome_referencia_figura1}
\end{figure}

\subsection{Modelo de Tabela}

A Tabela~\ref{tab:nome_tabela1} mostra um exemplo de tabela.

\begin{table}[h!]
\caption{Análise do IC nas linhas (L) e entrelinhas (E) de cana nas diferentes profundidades amostradas pelo índice de cone.}
\vspace{.2cm}
\centering

\begin{tabular}{ccccccccc}
\hline
Profundidade (m) & \multicolumn{2}{c}{0 a 0,1} & \multicolumn{2}{c}{0,1 a 0,2} & \multicolumn{2}{c}{0,2 a 0,3} & \multicolumn{2}{c}{0,3 a 0,4} \\ \hline
                 & L            & E            & L             & E             & L             & E             & L             & E             \\
Média (MPa)      & 1,39**       & 4,28**       & 1,86**        & 4,29**        & 2,20**        & 3,83**        & 2,46**        & 3,44**        \\
CV(\%)           & 54           & 57           & 55            & 54            & 46            & 49            & 48            & 43            \\ \hline
\end{tabular}
\\
\small{**:valores significativos para o nível de significância de 1\% pelo teste de Tukey; L – linhas; E – entrelinhas.}
\label{tab:nome_tabela1}
\end{table}

\section{Considerações Finais}

Demonstrar se os objetivos propostos foram alcançados, e as considerações finais da sua pesquisa (Helvetica, 9, Justificado).

\section*{Agradecimentos}

Campo destinado ao agradecimento à(s) agência(s) de fomento de financiamento do projeto de pesquisa. (Helvetica, 9, Justificado).


%Referências Bibliográficas
%\textbf{REFERÊNCIAS}

\bibliography{bibliografia.bib}
%\printbibliography[title={REFERÊNCIAS}]
%\printbibliography{bibliografia}



\end{document}
